\documentclass[12pt]{article}
\usepackage[margin=1in]{geometry}
\usepackage[T1]{fontenc}
\usepackage{lmodern}
\usepackage{parskip}

\begin{document}

\begin{center}
    \hfill Caleb Standfield\\
    \hfill UID: u1437560\\
    \hfill Games. Play. Japan.\\
    \hfill Professor Lazarus \\ \vspace{4mm}
    {\LARGE \textbf{Kemari Game Reflection}}\\[1em]
\end{center}

\vspace{1em}

The process of even coming up with the Kemari game was difficult; I spent a ton of time thinking about what I could do for my final project. A buddy of mine had made a cool rythm game in the past and I kinda wanted to do the same and Kemari stuck out as something I might be able to incorperate that form of rythm into.

That was my original goal, was a rythm game, and it wasn't untill after getting started did I realize that goal would be a little too ambitious, you would most definitely need some form of music and visuals to signal the player that they were hitting the rythm right. But then after thinking about it for awhile I realized that 1) I just didn't have the time or capactiy for that. And 2) most rythm games like Beat Saber, OSU, Geometry Dash, etc... all have some form of "winning". While it might have been possible to create the Kemari game with rythm game aspects and forgoing the whole concept of winning, I felt that the core of the game was more about elegance and team play, and thus I pivoted away from the rythm game aspect.

So what ended up being the hardest part about this project was getting the feel of the game down. (Even now I don't think its close to perfect) However, after much testing, trail and error, and play time, I think the core mechanics of building an elegance score and keeping the ball in the air as long as possible are complete.

For future students should they try something similar I would have a word of caution for them. The idea sounded simple at first, but once I got into development, there were a lot of problems. I had to deal with 3D models, textures, performance issues, and game logic all at once. Some of the assets I tried to use looked fine at first, but then caused huge problems in the engine, either because they were too detailed or because the textures and UVs were not set up in a way that worked well. I also ran into performance issues where the scene used way too much CPU and memory, which forced me to rethink how detailed the environment really needed to be. Lots of time sinks that took away from actually implementing the game.

As a student who has never done any form of game development, 3D model work, or anything like that, this was really hard to work with and definitely ate into my time alloted for this project. My scope seemed small but I had to spend a ton of time trying to work with these new tools that I quite literally had 0 experience with.

Another huge challenge was just the gameplay itself. I orginally thought it would be fairly simple, juggle the ball around, move a little bit to keep the ball under you, then pass to the next person. However, the logic behind having players moved around and still make passes accurate and keep the circle formation.... Well it got to be too much and I axed the movement and kept the players locked in place and tried to make the ball bath more determanistic. 

And just getting the NPCs to juggle the ball correctly killed me, they would let the ball touch the ground constantly or over shoot the other players, or fail to juggle the ball when receiving a pass. What really made the rythm game idea almost impossible is the fact that you have to play with NPCs, if this was a solo jugging game then I think it would have been more than possible to do that. But due to what Kemari is at its core, it just needed to be pivoted away from that type of game.

I did try to make certain parameters like juggling a certain amount of times before passing (but not too many) and make sure that you couldn't just spam buttons, so there is some rythm required to build elegance, but I wouldn't call it a rythm game at heart.

And I think that was the most important part of this process. Not trying to make Kemari into something that it isn't, at its core I understand Kemari to be about cooperation and thats what I tried to make without letting the scope of the project become too large. 

So at the end of all of this, I am glad I chose this project. It was definitely more difficult than I expected, I kinda thought I would be able to slap this game together in less that 80 hours, but at this point I've lost track of the time I have spent on this. I've definitely learned a lot about game dev and limiting scope (even if I did go over my time estimate), But more so than any of that, it made me think differently about both games and traditional culture, and it pushed me to solve problems in a more practical way. Even if the final version is not perfect, I think it reflects a real effort to take something historical and turn it into an interactive experience in a way that still respects what made it interesting in the first place, not trying to modernaize it to fit a trend.

\end{document}